\section{EDA Ringan/\textit{Compact}}

Sebagaimana telah diuraikan, EDA kanonik masih memerlukan penyimpanan populasi penuh pada setiap generasi.
Keterbatasan ini dijawab oleh algoritma \textit{compact}.
\textcite{khalfi_metaheuristics_2023} mendefinisikan algoritma \textit{compact} sebagai bentuk sederhana EDA yang menggantikan populasi eksplisit dengan satu \textit{Probability Vector} (PV)\label{first:pv} yang merepresentasikan populasi virtual berukuran $NP$, tanpa menyimpan seluruh individunya di memori.
Relevansi pendekatan ini ditunjukkan melalui contoh pada mikrokontroler ATmega328P dengan RAAM 2 KB, di mana algoritma genetika standar dengan 128 individu ber-\textit{encoding} 32-bit dapat menghabiskan sekitar 80\% memori, sehingga hampir tidak menyisakan ruang bagi proses lain \parencite{khalfi_metaheuristics_2023}.
Kondisi ini menggambarkan secara langsung tantangan memori pada perangkat \textit{edge-IoT}.

Secara umum, algoritma \textit{compact} diawali dengan inisialisasi PV dan \textit{sampling} solusi elit awal.
Pada setiap iterasi, satu solusi kandidat baru di-\textit{sampling} dari PV, lalu dibandingkan dengan solusi elit melalui skema turnamen pemenang-pecundang.
Selanjutnya, PV diperbarui berdasarkan perbedaan nilai variabel antara pemenang dan pecundang, dengan besar pembaruan yang berbanding terbalik terhadap ukuran populasi virtual (NP).
Solusi elit kemudian diganti apabila kandidat baru lebih baik \parencite{khalfi_metaheuristics_2023}.

Dengan mekanisme ini, hanya PV serta dua solusi, yaitu elit dan kandidat, yang perlu disimpan pada satu waktu.
Akibatnya, kebutuhan memori turun dari $O(NP \times D)$\label{first:symbol-onpd} pada EDA kanonik atau algoritma \textit{population-based} menjadi $O(D)$\label{first:symbol-od}, dengan $D$ menyatakan jumlah variabel keputusan.
Siklus kerja algoritma \textit{compact}, termasuk pembaruan PV dan keunggulan memorinya, diilustrasikan pada Gambar \ref{fig:siklus-compact}.

\begin{figure}[H]
    \centering
    \includegraphics[width=0.9\textwidth]{gambar_compact.png}
    \caption{Siklus Kerja Algoritma \textit{Compact} dan Mekanisme Pembaruan PV}
    \label{fig:siklus-compact}
\end{figure}

Gambar \ref{fig:siklus-compact} menunjukkan bahwa algoritma \textit{compact} menjalankan proses iteratif dengan hanya menyimpan satu PV serta dua solusi, sehingga kebutuhan memorinya tetap $O(D)$ tanpa menyimpan populasi virtual berukuran $NP$ secara eksplisit.
Prinsip pembaruan hemat memori ini diadopsi sebagai komponen \textit{compact} pada solusi yang diusulkan.

Pada ruang pencarian biner, \textit{compact-GA} merepresentasikan PV sebagai probabilitas per-bit yang diperbarui sebesar $\pm 1/NP$ berdasarkan perbedaan bit antara solusi pemenang dan pecundang \parencite{khalfi_metaheuristics_2023}.
Representasi ini sesuai untuk variabel diskret seperti yang terdapat pada ruang kebijakan ABAC.
Pengembangannya mencakup varian elitis serta \textit{Uniform Compact Genetic Algorithm} (UcGA)\label{first:ucga}, yang memungkinkan ukuran populasi virtual berbeda pada setiap variabel.

Untuk ruang kontine, \textit{Real-Valued Compact Genetic Algorithm} (rcGA)\label{first:rcga} menggunakan parameter distribusi Gaussian pada setiap dimensi, tetapi tetap menerapkan mekanisme pembaruan pemenang-pecundang.
Prinsip \textit{compact} juga diterapkan pada metaheuristik lain, seperti \textit{compact-PSO} yang mengadaptasi mekanisme pencarian PSO ke dalam representasi probabilistik \parencite{khalfi_metaheuristics_2023}.
Hal ini menunjukkan bahwa penggantian populasi eksplisit dengan model probabilistik ringan dapat diterapkan pada berbagai mekanisme pencarian.

Namun, algoritma \textit{compact} pada umumnya melakukan pencarian tak berkorelasi dan tidak menyimpan korelasi silang antar-variabel keputusan \parencite{khalfi_metaheuristics_2023}.
Dengan demikian, \textit{compact-GA}, rcGA, dan \textit{compact-PSO} standar pada dasarnya memodelkan setiap variabel secara independen.
Meskipun efisien dari sisi memori, asumsi ini membatasi kemampuan algoritma dalam merepresentasikan dependensi antar-atribut pada ruang kebijakan ABAC.
Kesenjangan antara efisiensi memori dan kemampuan memodelkan dependensi tersebut menjadi salah satu motivasi utama penelitian ini.